% Bagian: computability
% Bab: computability-theory
% Subbagian: universal-part-function

\documentclass[../../../include/open-logic-section]{subfiles}

\begin{document}

\olfileid[id]{cmp}{thy}{uni}
\olsection{Fungsi Dapat Dihitung Parsial Universal}

\begin{thm}
\ollabel{thm:univ-comp}
Terdapat suatu fungsi dapat dihitung parsial universal~$\fn{Un}(e,x)$. Dengan
kata lain, terdapat suatu fungsi $\fn{Un}(e,x)$ sedemikian sehingga:
\begin{enumerate}
\item $\fn{Un}(e,x)$ dapat dihitung parsial.
\item Jika $f(x)$ merupakan sembarang fungsi dapat dihitung parsial, maka
  terdapat suatu bilangan asli $e$ sedemikian sehingga
  $f(x) \simeq \fn{Un}(e,x)$ bagi setiap~$x$.
\end{enumerate}
\end{thm}

\begin{proof}
Tetapkan $\fn{Un}(e,x) \simeq U(\umin{s}{T(e,x,s)})$, dengan $U$ dan $T$
seperti dalam teorema bentuk normal Kleene
(\olref[nfm]{thm:normal-form}).
\end{proof}

\begin{explain}
Ini hanyalah cara yang tepat untuk mengatakan bahwa kita mempunyai enumerasi
efektif atas fungsi-fungsi dapat dihitung parsial. Gagasannya ialah: jika kita
menulis $f_e$ untuk fungsi yang didefinisikan oleh
$f_e(x) = \fn{Un}(e,x)$, maka urutan $f_0$, $f_1$, $f_2$, \dots mencakup semua
fungsi dapat dihitung parsial, dengan sifat bahwa $f_e(x)$ dapat dihitung
``secara seragam'' dalam $e$ dan~$x$. Demi kesederhanaan, kita menggunakan
fungsi biner yang universal bagi fungsi berargumen satu, tetapi dengan
mengodekan barisan bilangan kita dapat dengan mudah memperumumnya untuk lebih
banyak argumen. Misalnya, perhatikan bahwa jika $f(x,y,z)$ merupakan fungsi
rekursif parsial berargumen~$3$, maka fungsi
$g(x) \simeq f((x)_0, (x)_1, (x)_2)$ merupakan fungsi rekursif parsial
berargumen satu.
\end{explain}

\end{document}
